\chapter{Orthodromic Rivers: Water and Governance}
\label{ch:18}

The geometry of the preceding chapter determines where a freshwater corridor might run. The hydraulics determines what it would cost to run water through it, and the governance determines whether it could be built and kept in operation. This chapter treats the second and third. The hydraulics is dominated by two relations that every proposal for long-distance conveyance must satisfy, the energy required to lift water and the energy dissipated by friction, and the chapter shows that the second, which is less discussed, is the more severe constraint over the distances that the concept contemplates. The governance is dominated by the fact that the water of a continent is claimed by many parties, and that the claims are not settled by engineering. The assessment concludes that the concept is in the second readiness class in its hydraulic aspects and, as to governance, depends on arrangements whose feasibility is open to doubt. Baseline data on scarcity, withdrawal, and transboundary basins are taken from the standing international assessments~\cite{unesco2024,fao2024}, and the numerical parameters of this chapter should be read as stipulated illustrations against that background.

Begin with lift. The work required to raise a cubic meter of water through a height of one kilometer is about $9.8\times10^{6}$ joules, which is $2.7$ kWh, and it is independent of the length of the route over which the rise is accomplished. This figure is not a design parameter but a floor, and it supplies a standard of comparison. Reverse-osmosis desalination of seawater at present consumes about three to four kilowatt-hours per cubic meter in large plants, including pretreatment and intake, and the thermodynamic minimum for seawater of typical salinity is a little over one kilowatt-hour per cubic meter. A corridor that lifted water by a kilometer would therefore consume about as much energy per unit volume as a desalination plant that supplied the same water at the coast, before any allowance for friction, and a corridor that lifted it by more than a kilometer would consume more. For a conduit with a uniform mean uphill gradient $S$, so that the rise is proportional to the distance traveled, there is a distance at which the energy of lift equals the energy of desalination, and the formalization computes it as about $10{,}900$ km for $S=10^{-4}$ and $3.5$ kWh per cubic meter at a pumping efficiency of $0.85$. Within that distance, transport by uphill gradient is cheaper in energy than desalination at the coast; beyond it, it is dearer. The comparison is limited to energy and ignores capital cost, which is the larger part of the cost of either, but it shows that the distances at issue are not so great as to make lift prohibitive when the gradient is gentle.

Friction, however, is another matter. Water moving through a conduit loses head in proportion to the conduit's length and to the square of its velocity, and inversely to its diameter. The loss must be made up by pumping when the conduit does not descend as fast as the friction demands, and it must be made up by slope when the conduit is gravity-fed. In the first case it is a continuing energy cost, and in the second it is an elevation that must be available. The formalization shows that, for a given discharge $Q$, the friction slope varies as $Q^2/D^5$, which means that the diameter is the dominant design variable: doubling it reduces the friction by a factor of thirty-two. A pipeline of ten meters in diameter, carrying about $157\ \mathrm{m^3/s}$ at $2$ m/s, with a friction factor $0.012$, loses $1{,}223$ meters of head over $5{,}000$ km, which costs $3.9$ kWh per cubic meter at $85$ per cent efficiency, more than the whole cost of desalination. Widening the pipe to twelve meters reduces the figure to $1.6$ kWh, to fourteen meters to $0.73$ kWh, and to sixteen meters to $0.37$ kWh. The sensitivity illustrates why the concept cannot be assessed by a single number and why it is a problem of design in which the capital cost of larger conduits is balanced against the energy cost of friction. Open channels dissipate energy by a different law, expressed by the Manning equation~\cite{manning1891}, in which the velocity is proportional to the square root of the slope. A concrete-lined channel one hundred meters wide and ten meters deep, on a gradient of $10^{-4}$, would carry about $3{,}400\ \mathrm{m^3/s}$, or about $108$ cubic kilometers a year, and would fall $500$ meters in $5{,}000$ km, so that a gravity-fed channel of such length must descend from a source at least that much above the point of delivery. The geographical requirement is a severe one; the sources of fresh water on land are not usually higher than the regions short of water by the necessary amount, and the opposite relation is common.

Open channels expose the water to the atmosphere and therefore to evaporation and, in the absence of a lining, to seepage. Evaporation from the channel just described, in a region where the open-water evaporation rate is two meters per year, would amount to about one cubic kilometer a year over $5{,}000$ km of length, or about nine-tenths of one per cent of the flow. This loss is not a serious one for a channel of such capacity, but it increases in proportion to the surface and decreases in relation to the discharge, so that small channels over long distances are lossy and large channels are not. Seepage depends on the lining and on the substrate, and a failure of the lining over an arid region may deliver the water to the aquifer and not to the destination, which may be a benefit in a recharge program and a defect in a delivery program. These considerations point to a conclusion that is not new but is worth repeating: the efficient scale of a long conduit is large, and large works concentrate risk, so that the cost of failure and the difficulty of repair are also large.

Ecological effects are the third class of physical consequence. A channel that connects drainage basins previously separate creates a path for the movement of organisms, and the history of canals contains instructive cases. The Suez Canal opened a corridor through which species of the Red Sea have invaded the Mediterranean, with substantial effects on its fisheries and ecology, and proposals to separate the Great Lakes from the Mississippi system are motivated by the fear that species of carp will cross between them through the Chicago waterway. A freshwater corridor of continental length would connect basins that have been isolated for millennia, and the movement of pathogens, plants, and animals through it would have to be prevented by design, which is difficult, or accepted as a cost, which may be unacceptable. The corridor also intercepts the flow of surface and subsurface water across its path and may disrupt wetlands, migrations, and sediment transport. An assessment of these effects is a prerequisite for any decision to proceed and could not be left to the later stages of design.

Governance determines more than the physical aspects what can be done. Water is claimed by those who live on a river and by those who live far from it, by the upstream state and by the downstream state, by farms and cities and the ecosystems on which they depend, and the law of transboundary waters has developed principles, notably equitable and reasonable utilization and the obligation not to cause significant harm to other states, that are widely accepted and unevenly applied. The Indus Waters Treaty of 1960 has endured conflict between the parties and is often cited as a rare success, and the failure of the parties on the Nile to agree on the allocation of its waters is often cited as the contrary case. A corridor that transfers water between basins is a physical proposition that cannot be separated from these claims. It would alter the allocation among the parties, and the parties' consent would have to be obtained for a project whose benefits and burdens fall on different persons. The design principles of Ostrom for durable common-pool arrangements imply that the proposal should provide for monitoring by the users, for graduated sanctions for violation of the rules, for accessible mechanisms for the resolution of disputes, and for the recognition by higher authorities of the right of the users to organize. A project designed by a central body and imposed on the users, however elegant its geometry, would fail these tests.

The assessment under the standard of Chapter 16 follows. The mechanism is explicit and the energy balance has been written out. The relation of lift to the alternatives has been stated, and the friction analysis shows that the designs for which the energy cost is competitive with desalination require conduit dimensions much larger than any in service, which is the main reason that the concept is placed in the second readiness class and not the first. Dependencies include the production of materials in the required quantity, the availability of sources at sufficient elevation, and the agreement of the states concerned. The comparison with alternatives that do not require the construction of the corridor, including demand management, wastewater reuse, local recharge, and desalination, is mandatory, and for the present these are the measures on which present obligations to supply water may legitimately be based. The corridor remains a candidate for study as a supplement to them and not as a substitute.

\section*{Formalization}

Let water of density $\rho=1000\ \mathrm{kg\,m^{-3}}$ be moved at discharge $Q$ through a conduit of length $L$, diameter $D$, friction factor $f$, with net rise $\Delta z$ and pumping efficiency $\eta\in(0,1]$. The friction slope is $S_f=f\,v^2/(2gD)$ with $v=4Q/(\pi D^2)$, and the energy required per unit volume, ignoring minor losses, is
\[
\frac{E}{V}=\frac{\rho g\,(\Delta z+S_fL)}{\eta}.
\]

\begin{proposition}
(i) At fixed $Q$, $S_f=8fQ^2/(\pi^2gD^5)$, so that $S_f$ is proportional to $D^{-5}$ and a factor $\lambda>1$ in diameter reduces friction by $\lambda^{-5}$. (ii) For a uniform mean uphill gradient $S$, so that $\Delta z=SL$, and neglecting friction, the energy of lift equals a reference energy $E_d$ per unit volume (for example, that of desalination) at the break-even distance
\[
L^\ast=\frac{\eta\,E_d}{\rho g\,S},
\]
and lift is energetically cheaper than the reference for $L<L^\ast$. (iii) If annual cost per unit length is $a\,D^\beta+b\,Q^3D^{-5}$, where the first term is capital cost with $\beta>0$ and the second is the cost of the energy to overcome friction, then the minimizing diameter is $D^\ast=(5bQ^3/(\beta a))^{1/(\beta+5)}$, and at the optimum the energy term equals $\beta/5$ times the capital term.
\end{proposition}

\begin{proof}
(i) Substituting $v=4Q/(\pi D^2)$ into $S_f$ gives $S_f=f\,16Q^2/(2g\pi^2D^5)$. (ii) Setting $\rho gSL^\ast/\eta=E_d$ and solving. (iii) The pumping power to overcome friction is $\rho gQS_fL/\eta\propto Q^3D^{-5}$, so the cost per unit length has the stated form. Setting the derivative $\beta aD^{\beta-1}-5bQ^3D^{-6}$ to zero gives $D^{\beta+5}=5bQ^3/(\beta a)$ and at that point $bQ^3D^{-5}=(\beta/5)aD^\beta$.
\end{proof}

Numerical values follow. With $E_d=3.5\ \mathrm{kWh\,m^{-3}}=1.26\times10^{7}$ J m$^{-3}$ and $S=10^{-4}$: $L^\ast=1.2844\times10^7$ m $\approx12{,}844$ km at $\eta=1$ and $\approx10{,}917$ km at $\eta=0.85$. The lift energy for $\Delta z=1000$ m is $\rho g\Delta z=9.81\times10^6$ J m$^{-3}=2.725$ kWh m$^{-3}$ at $\eta=1$. For $Q=157\ \mathrm{m^3/s}$, $f=0.012$, $L=5000$ km and $\eta=0.85$, with $\Delta z=0$ (a level conduit), the friction energy is $3.92$ kWh m$^{-3}$ at $D=10$ m, $1.58$ at $12$ m, $0.73$ at $14$ m, and $0.37$ at $16$ m. The friction factor is held constant for simplicity; in practice it varies weakly with the Reynolds number and the relative roughness, and this variation alters the figures by tens of per cent but not their scaling with $D$. For the open channel, the Manning equation $v=n^{-1}R_h^{2/3}S^{1/2}$, with $n=0.012$, $R_h=1000/120=8.33$ m and $S=10^{-4}$, gives $v=3.43$ m/s and $Q=3425\ \mathrm{m^3/s}$, which is $108\ \mathrm{km^3\,yr^{-1}}$, and the evaporative loss over $5000$ km at $2$ m yr$^{-1}$ and width $100$ m is $1.0\ \mathrm{km^3\,yr^{-1}}$, a fraction $0.93$ per cent. The rule of (iii), which parallels the classical law for the choice of the cross-section of an electrical conductor, is useful in practice: at the optimal design the friction energy is a fixed fraction of the capital cost, so that if the diameter is chosen optimally and the cost of energy is higher than the proposal assumes, the design should be revised toward wider conduits and not narrowed.
